\lecture{4}{Thu 30 Mar 2023 09:27}{Volume of Parallelopiped}

A \textit{parallelopiped} is of the form
\[
	L([0,1]^k)
.\] 
where $L: \mathbb{R}^k \to \mathbb{R}^n$, $k\le n$, is a linear map. We say $L e_i = x_i$. The parallelopiped is spanned by $x_1, \ldots, x_k$.

The desired properties of the volume $V(X)$ are:
\begin{itemize}
	\item (Rotation Invariance) If $A^T A = I$, then $V(AX) = V(X)$.
	\item (Projection Invariance) If $X$ is entirely contained in $\mathbb{R}^k \times \{0\}$, such that $x_i = z_i \oplus 0_{\mathbb{R}^{n-k}}$, then the volume $V(X)$ concides with the $k$-dimensional parallelopiped.
\end{itemize}

\begin{theorem}
	\[
		V(X) = \text{det}(X^T X)^{\frac{1}{2}}
	.\] 
	satisfies the desired properties.
\end{theorem}
\begin{proof}
	Let $F(X) = \text{det }(X^T X)$, so 
	\[
		F(AX) = \text{det}\left( (AX)^T (AX) \right) = \text{det}(X^T A^T A X) = \text{det} (X^T X) = F(X)
	.\] 

	Also, if $x_k = z_k \oplus 0$, then
	\[
		X^T X = 
	\begin{bmatrix}
		Z^T & 0 \\
	\end{bmatrix}
	\begin{bmatrix}
		Z \\
		0 \\
	\end{bmatrix}
	= Z^T Z
	.\] 

	Hence 
	\[
		f\left( X \right)  = \text{det}(Z^T Z) = \text{det}Z^T \text{det}Z =( \text{det}Z)^2
	.\] 

\end{proof}



\begin{lemma}
	$X^T X$ is symmetric nonnegative matrix (positive semidefinite)
\end{lemma}
\begin{proof}
	If $Q = X^T X$, we show $u \cdot Q u \ge 0$.
	$u \cdot X^T Xu = Xu \cdot Xu \ge 0$.

	Note
	\[
		(X^T X)^T = X^T X
	.\] 
	Hence it is symmetric, ad with nonnegative eigenvalues.
\end{proof}

Therefore 
	$\text{det}(X^T X) = \lambda_1 \ldots \lambda_k\ge 0$. Thus $F(X) \ge 0$, $V(X) = F(X)^{\frac{1}{2}}$ is well-defined. 


\begin{lemma}
	If $W$ is a $k$-dimensional subspace of $\mathbb{R}^n$, then there is an orthonormal basis $\{b_1, \ldots, b_n\} $ of $\mathbb{R}^n$ such that $b_1, \ldots, b_k$ is a basis for $W$.
\end{lemma}
\begin{proof}
	Apply the Gram-Schmidt process.
\end{proof}


We claim:
\begin{proposition}
	For any $x_1, \ldots, x_k$, there is a rotation $h(x) = Ax$ such that  $A x_1, \ldots, Ax_k \in \mathbb{R}^k \times \{0\} $.
\end{proposition}
\begin{proof}
	Let $W = \left\langle x_1, \ldots, x_k \right\rangle $ subspace spanned by $x_1, \ldots, x_k$. WLOG assume $\operatorname{dim} W = k$. Then, let $\{b_1, \ldots, b_n\} $ be orthonormal basis be as in lemma. Define $h$ by
	\begin{align*}
		h: \mathbb{R}^n &\longrightarrow \mathbb{R}^n \\
		b_i &\longmapsto h(b_i) = Ab_i = e_i 
	.\end{align*}
	Thus $AB = I$, $B = [b_1, \ldots, b_n] $.
	\[
		A(W) \subset \mathbb{R}^k \times \{0\} 
	.\] 
\end{proof}

Once we have the claim, we have uniqueness of the volume:
\[
	V(X) = V(AX) = V([Z;0]) = |\text{det}Z| = \text{det}(X^T X)^{\frac{1}{2}}
.\] 


\begin{theorem}[Pythagorean Theorem]
	\[
		V(X)^2 = \sum_{[I]} \text{det}(X_I)^2
	\]
	where $[I]$ is sum over asending $k-$ tuples.
\end{theorem}
\begin{example}

\begin{itemize}
\item $k=1$.
	
	\begin{align*}
		x &= (x_1, \ldots, x_k)\\
		\|x\|^2 &=  x_1^2 + \ldots + x_k^2 
	.\end{align*}
	This coincides with the elementary Pythagorean theorem.
\item $k = 2, n = 3$
\[
	X = [X_1, X_2] = 
 	\begin{bmatrix}
		x_{11} & x_{21} \\
		x_{12} & x_{22} \\
		x_{13} & x_{23} \\
	\end{bmatrix}
\]
Hence
\[
	V(X)^2 = 
\begin{vmatrix}
	x_{11} & x_{21} \\
	x_{12} & x_{22} \\
\end{vmatrix}
+ \ldots \text{ (3 minors in total)}
= |X_1 \times X_2|^2
.\] 

\item $k = n$, $X_I = X$.

	Now $V(X)^2 = \text{det}(X)^2 = \text{det}(X^T X)$
\end{itemize}
\end{example}

\begin{proof}
	Cases $k = 1$ and $k=n$ are immediate.

	Let $F(X) = \text{det}(X^T X)$ and $G(X) = \sum_{[I]}\text{det}(X_I)^2$.

	The elementary row operations on $M$ is $EM$. The column operation can be done via transpose:
	\[
		(EM^T)^T = ME^T
	.\] 

	Let $E$ be elementary row operations, then
	\begin{align*}
		F(X E^T) &=  \text{det}\left( (X E^T) ^T (XE^T) \right)  \\
		&= \text{det}E \text{det}X^T X \text{det }E^T \\
		&= F(X) \text{det}\left(E \right) ^2\\
		&= F(X) 
	.\end{align*}
	Note that
	\[
		(XE^T)_I = X_I E^T
	.\] 
	Hence
	\begin{align*}
		G(X E^T) &= \sum_{[I]} \text{det}(X_I E^T)^2\\
			 &= \sum_{[I]}\text{det}(X_I)^2 \text{det}(E)^2 \\
			 &= \sum_{[I ]}\text{det}(X_I)^2 \\
			 &= G(X) 
	.\end{align*}

	\textbf{Claim}. We can transform $X$ to the form
	\[
	\begin{bmatrix}
		b_1 & b_2 & \ldots & b_{k-1} & b_k \\
		0 & 0 & \ldots & 0 & \lambda \\
	\end{bmatrix}
\]
where $\lambda\ge 0$ and $b_1, \ldots, b_k \in \mathbb{R}^{n-1}$ and orthogonal.

\end{proof}



